% SPDX-License-Identifier: Apache-2.0
% Sections 2--4 adapt the Apache-2.0 source of OpenAI's 2026 pi manuscript.
% Attribution and the original Apache-2.0 license are preserved in NOTICE and LICENSE.
\documentclass[11pt]{article}
\usepackage[T1]{fontenc}
\usepackage{lmodern,fancyhdr}
\usepackage[margin=1in]{geometry}
\usepackage{amsmath,amssymb,amsthm,mathtools,booktabs,longtable}
\usepackage[colorlinks=true,linkcolor=blue,urlcolor=blue,citecolor=blue]{hyperref}
\setlength{\emergencystretch}{3em}
\newcommand{\CC}{\mathbb C}
\newcommand{\RR}{\mathbb R}
\newcommand{\QQ}{\mathbb Q}
\newcommand{\ZZ}{\mathbb Z}
\newcommand{\NN}{\mathbb Z_{\ge0}}
\newcommand{\ii}{\mathrm i}
\DeclareMathOperator{\ord}{ord}
\DeclareMathOperator{\lcm}{lcm}
\DeclareMathOperator{\codim}{codim}
\DeclareMathOperator{\Spec}{Spec}
\DeclareMathOperator{\Proj}{Proj}
\DeclareMathOperator{\Supp}{Supp}
\newtheorem{theorem}{Theorem}[section]
\newtheorem{proposition}[theorem]{Proposition}
\newtheorem{lemma}[theorem]{Lemma}
\newtheorem{corollary}[theorem]{Corollary}
\theoremstyle{definition}
\newtheorem{definition}[theorem]{Definition}
\theoremstyle{remark}
\newtheorem{remark}[theorem]{Remark}
\numberwithin{equation}{section}

\pagestyle{fancy}
\fancyhf{}
\fancyhead[L]{\small Candidate logarithm argument}
\fancyhead[R]{\small Not independently certified}
\fancyfoot[C]{\thepage}
\setlength{\headheight}{14pt}
\hypersetup{pdftitle={Irrationality exponents of log 2 and rational logarithms: candidate proof},
pdfauthor={Research draft prepared with Codex},
pdfsubject={Proposed distinct-center interpolation argument; unverified mathematical status}}
\title{Candidate proofs for the irrationality exponents\\of $\log 2$ and rational logarithms}
\author{Research draft prepared with Codex}
\date{8 October 2026}
\begin{document}
\maketitle
\begin{abstract}
This document gives the full proposed interpolation and determinant
argument for the irrationality exponent of the natural logarithm of two.
It adapts the separated-weight method in OpenAI's manuscript on $\pi$,
replacing centers with $Y=1$ by distinct multiplicative centers $Y=2^j$.
A separate section extends the argument to $\log a$ for positive rational $a\ne1$,
nonzero rational linear combinations of such logarithms, and positive
algebraic numbers with a positive integral power in $\QQ$.
The crucial generalized interpolation statement and the resulting proof
are candidates requiring further verification. This document does not
establish a certified new theorem.
\end{abstract}
\begin{center}
\fbox{\begin{minipage}{0.91\textwidth}
\textbf{Verification status.}
The new distinct-$Y$ geometric interpolation argument has been read
internally, but has not been independently certified or formalized.
Parameter calculations and successful TeX compilation do not validate
that argument. Every proof below is presented as part of the candidate
derivation, with its exact claim available for scrutiny.
\end{minipage}}
\end{center}
\tableofcontents
\clearpage

\section{Claim, provenance, and verification boundary}
For an irrational real number $x$, use
\[
\mu(x)=\sup\{\nu>0:0<|x-p/q|<q^{-\nu}
\text{ for infinitely many coprime }p\in\ZZ,\ q\ge2\}.
\]
Dirichlet's theorem gives $\mu(x)\ge2$. The published bound used here
as background is Marcovecchio's $\mu(\log2)<3.57455391$
\cite[Theorem 1.1]{marcovecchio}.

Sections 2--4 adapt OpenAI's \emph{The irrationality exponent of $\pi$
is $2$}, dated 24 September 2026 \cite{openai}. Credit for the local
multiplicity, projective interpolation, determinant and parameter
mechanisms belongs to that manuscript and the earlier work cited below.
The original source is publicly linked in the references;
the Apache-2.0 license and adaptation notice accompany this manuscript. This standalone source has no external TeX inputs,
figures, or bibliography files.

The substantive proposed modification is to use centers
$(2^j,jp_1/q_1,\ldots,jp_m/q_m)$ and local coordinates
\[
Y=2^j(1+t),\qquad
X_i=jp_i/q_i+u_i+\log(1+t).
\]
The new geometric statement must allow successive weight thresholds to
be selected independently of the center coordinates. Only the eventual
degree threshold may depend on the fixed centers. This order of
quantifiers is essential to the final contradiction.

The posted Lean implementation for the source $\pi$ theorem has not been
rebuilt or exhaustively audited here. It does not certify this proposed
$\log2$ extension. No exhaustive novelty claim is made. An external
mathematical audit or an appropriate formal proof remains necessary to
establish the candidate result.

\begin{theorem}[Proposed result; candidate proof]
\label{thm:main}
For every real $\nu>2$, there is an integer $Q(\nu)$ such that
\[
|\log2-p/q|\ge q^{-\nu}
\qquad(p\in\ZZ,\ q\in\ZZ,\ q\ge Q(\nu)).
\]
Consequently $\mu(\log2)=2$.
\end{theorem}
The remaining sections give the proposed argument, rather than a claim
that its correctness has already been established.

\section{Interpolation with separated weights}
\label{geom:section}

We prove the interpolation statement used in the determinant argument. Its
essential uniformity is that the weights can be chosen before the interpolation
points: each coordinate weight has a lower threshold depending on the earlier
weights, but not on the locations of the points. The degree at which
interpolation becomes surjective is allowed to depend on all the fixed data.

All multi-indices below have nonnegative integer entries. For a positive
rational vector $W=(w_0,\ldots,w_m)$, the $W$-degree of the
monomial $Y^hX_1^{\alpha_1}\cdots X_m^{\alpha_m}$ is
$w_0h+\sum_iw_i\alpha_i$. Write $\mathcal P_W(H)$ for the complex vector
space of polynomials of $W$-degree at most $H$. For a positive rational vector
$V=(v_0,\ldots,v_m)$, let $\mathcal J_V(H)$ be the quotient of
$\CC[[t,u_1,\ldots,u_m]]$ by the ideal of series whose monomials all have
$V$-weight at least $H$. Thus a vector in $\mathcal J_V(H)$ is a packet of
coefficients indexed by
\[
 v_0s+\sum_{i=1}^m v_i\beta_i<H.
\]
The strict inequality in this definition will be used throughout.

\begin{theorem}[Separated-weight interpolation]
\label{geom:interpolation}
Fix integers $m,K\ge1$ and positive rational numbers $w_0,v_0,\theta$ such
that $0<\theta<1$ and
\begin{equation}
\label{geom:volume-hypotheses}
 K\theta^m<1,
 \qquad K\frac{w_0}{v_0}\theta^m<1.
\end{equation}
There are successive lower thresholds for positive rational numbers
$w_1,\ldots,w_m$, depending only on these fixed data and the previously
chosen weights, with the following property. Put
\[
 W=(w_0,w_1,\ldots,w_m),\qquad
 V=(v_0,w_1/\theta,\ldots,w_m/\theta).
\]
For any points
\[
 a_j=(y_j,c_{j1},\ldots,c_{jm})\in\CC^{m+1},\qquad 0\le j<K,
\]
where $y_0,\ldots,y_{K-1}$ are nonzero and pairwise distinct, the map
\begin{equation}
\label{geom:jet-map}
 \begin{aligned}
 \mathcal P_W(H)&\longrightarrow\bigoplus_{j=0}^{K-1}\mathcal J_V(H),\\
 P&\longmapsto
 \left(P\bigl(y_j(1+t),(c_{ji}+u_i+\log(1+t))_{i=1}^m\bigr)\right)_{j=0}^{K-1}.
 \end{aligned}
\end{equation}
is surjective for all sufficiently large, sufficiently divisible integers $H$.
The divisibility requirement can be chosen using only the weights. The
remaining lower threshold for $H$ may depend on the points.
\end{theorem}

We first establish a numerical inequality on algebraic curves, and then use
it to obtain all the jets in \eqref{geom:jet-map}. The local estimate driving
the curve argument is recorded separately.

\subsection{A local multiplicity comparison}

A collection of vector fields is called a \emph{normal basis} to a smooth
subvariety at a point if its images form a basis of the ambient tangent space
modulo the tangent space of the subvariety. For an irreducible subvariety,
this terminology will refer to a general smooth point. There are only
finitely many possible index sets in the applications below, so all their
generic rank conditions can be checked on one common open subset.

\begin{lemma}[Weighted transverse multiplicity]
\label{geom:multiplicity}
Give the coordinates $x_1,\ldots,x_d$ positive rational degree weights
$\rho_1,\ldots,\rho_d$. On an open subset of $\CC^d$, let
$D_1,\ldots,D_d$ be commuting analytic vector fields forming a frame, and
assign them positive rational costs $\kappa_1,\ldots,\kappa_d$. Fix
$\varepsilon>0$.

Let $N>0$. Suppose that polynomials $f_1,\ldots,f_r$ have weighted degree at most $N$,
that $Z$ is an irreducible component of their common zero set of codimension
$k\ge1$, and that
\begin{equation}
\label{geom:persistent-derivatives}
 D^\gamma f_\ell\big|_Z=0
 \quad\text{whenever}\quad
 \sum_b\kappa_b\gamma_b<\varepsilon N.
\end{equation}
For any coordinate normal basis $A$ and any $D$-field normal basis $B$,
each of cardinality $k$, one has
\begin{equation}
\label{geom:multiplicity-comparison}
 \prod_{a\in A}\rho_a
 \le 2k!\,\varepsilon^{-k}\prod_{b\in B}\kappa_b
\end{equation}
The bound is uniform in the equations, the component $Z$, and the
analytic frame; no lower threshold for $N$ is required.
\end{lemma}

\begin{proof}
Choose a smooth point $x$ of $Z$ where both normal-basis conditions hold and
which lies on no other irreducible component of the common zero set. Flowing
from $Z$ along the fields indexed by $B$ gives local analytic coordinates
$(z,b)$, where $z$ parametrizes $Z$ and $b=(b_1,\ldots,b_k)$ consists of
transverse flow parameters. Commutativity implies that the chosen fields
are the derivatives in the $b$-directions. Thus
\eqref{geom:persistent-derivatives} says that in the expansion
\[
 f_\ell(z,b)=\sum_{\alpha\in\NN^k} f_{\ell,\alpha}(z)b^\alpha,
\]
every coefficient function of $B$-weight less than $\varepsilon N$ vanishes
identically on the chosen neighborhood in $Z$.

Let $S$ be the affine coordinate slice through $x$ obtained by fixing the
coordinates outside $A$. It is transverse to $Z$, and the zero set of the
restricted equations is isolated at $x$. The flow parameters $b$ give
regular analytic coordinates on $S$; the coordinates $z$ become analytic
functions $z(b)$. Substitution cannot generate a missing term of smaller
$B$-weight. Indeed, each surviving summand already contains a monomial
$b^\alpha$ of weight at least $\varepsilon N$, and its analytic coefficient
has only nonnegative powers of $b$. Consequently the ideal of the
restricted equations is contained in the monomial ideal of weight at least
$\varepsilon N$. Its quotient therefore has length at least
\begin{equation}
\label{geom:local-length-lower}
 \#\left\{\alpha\in\NN^k:
       \sum_{b\in B}\kappa_b\alpha_b<\varepsilon N\right\}
 \ge \frac{(\varepsilon N)^k}{2k!\prod_{b\in B}\kappa_b}
\end{equation}
In fact the factor $1/2$ can be omitted: the simplex
$\{x\in\RR_{\ge0}^k:\sum_{b\in B}\kappa_bx_b<\varepsilon N\}$
is covered by the disjoint half-open unit cubes $\alpha+[0,1)^k$
indexed by the admissible multi-indices $\alpha=\lfloor x\rfloor$. Its volume is
$(\varepsilon N)^k/(k!\prod_{b\in B}\kappa_b)$. Thus the displayed
bound holds for every $N>0$. Analytic local coordinates preserve finite
length, so no radius or coefficient bound for the flow map is needed.

In the regular local ring of the $k$-dimensional slice, the ideal generated
by the restricted equations is primary to the maximal ideal. Over the
infinite field $\CC$, $k$ sufficiently general constant linear
combinations of these equations still have an isolated common zero at
$x$: successive generic combinations avoid the finitely many
positive-dimensional components that remain. Their ideal is contained
in the original ideal. Their quotient length is therefore an upper bound
for the length in \eqref{geom:local-length-lower}.

Each of the $k$ combinations has weighted degree at most $N$ in the
variables indexed by $A$. We claim that their local intersection length
is at most
\begin{equation}
\label{geom:weighted-bezout}
 \frac{N^k}{\prod_{a\in A}\rho_a}.
\end{equation}
Choose an integer $M_0$ clearing the denominators of the $\rho_a$, and
make the substitutions
\[
 x_a=\lambda_a+z_a^{M_0\rho_a}\qquad(a\in A).
\]
Choose the constants $\lambda_a$ so that no preimage of the given point
is ramified. There are $\prod_{a\in A}(M_0\rho_a)$ such preimages, each
with the same local intersection length. The substituted polynomials
have ordinary total degree at most $M_0N$. The ordinary B\'ezout bound
for isolated complete-intersection zeros, with their local
multiplicities, bounds the sum of these lengths by $(M_0N)^k$;
see \cite[Corollary~VIII.3, p.~144]{mondal2021}.
The multiplicities there are local quotient lengths, unchanged by
completion \cite[Section~IV.4.2]{mondal2021}.
Dividing gives \eqref{geom:weighted-bezout}.
This application concerns isolated zeros only and permits other
components elsewhere.

Combining \eqref{geom:local-length-lower} and
\eqref{geom:weighted-bezout} proves the result. Neither the
polynomials nor the possibly varying component enters the constant in
\eqref{geom:multiplicity-comparison}.
\end{proof}

\subsection{The curve inequality}

Choose a positive rational number $\sigma$ sufficiently small that
\begin{equation}
\label{geom:sigma}
 \begin{split}
 (1+3\sigma)^{m+1}K(w_0/v_0)\theta^m&<1,\\
 (1+3\sigma)^mK\theta^m&<1,\\
 (1+\sigma)\theta&<1.
 \end{split}
\end{equation}
This choice is possible by \eqref{geom:volume-hypotheses}. Set $v_i=w_i/\theta$
for $i>0$. Fix $L_0=m+2$ and put
\begin{equation}
\label{geom:constant}
 C_* =\max_{1\le k\le m}2k!(L_0/\sigma)^k.
\end{equation}
We choose the positive weights successively so that for every pair of
subsets $A,B\subseteq\{0,\ldots,m\}$ of the same cardinality between
$1$ and $m$,
\begin{equation}
\label{geom:separation}
 \prod_{a\in A}w_a>C_*\prod_{b\in B}v_b
\end{equation}
whenever the largest positive index in their symmetric difference belongs
to $A$.

Here and below a condition involving that largest index is imposed only
when a positive index does differ. The successive choice is possible:
if this largest index is $i$, the ratio in \eqref{geom:separation} is
$w_i$ times a positive expression in the already chosen lower weights
and the fixed zeroth weights. Every higher index occurs in both sets or
in neither; in the former case its ratio is $w_j/v_j=\theta$. Hence the
expression is independent of all higher weights. There are finitely many
pairs of sets. A sufficiently large choice of $w_i$ enforces all
conditions whose largest differing index is $i$, and subsequent choices
do not disturb them. This construction involves no interpolation point.

Let $C$ be an irreducible algebraic curve whose generic point lies in the
affine coordinate chart. Such a curve has a smooth projective model over
$\CC$; see \cite[Theorem~53.2.6 and Lemma~53.2.8, Tag~0BXX]{stacks}.
On that model define
\begin{equation}
\label{geom:curve-degree}
 \deg_W C=
 \sum_P\max\left\{0,-\frac{\operatorname{ord}_P Y}{w_0},
                  -\frac{\operatorname{ord}_P X_1}{w_1},\ldots,
                  -\frac{\operatorname{ord}_P X_m}{w_m}\right\}.
\end{equation}
The order of an identically zero coordinate is understood as $+\infty$;
it contributes no pole. At a branch $P$ through $a_j$, set
\begin{equation}
\label{geom:branch-contact}
 h_P=\min\left\{
 \frac{\operatorname{ord}_P(Y-y_j)}{v_0},
 \frac{\operatorname{ord}_P(X_1-c_{j1}-\log(Y/y_j))}{v_1},\ldots,
 \frac{\operatorname{ord}_P(X_m-c_{jm}-\log(Y/y_j))}{v_m}
 \right\}.
\end{equation}
The logarithm here is its formal expansion at $Y=y_j$, namely $\log(Y/y_j)$. All entries have
positive order, possibly infinite, and their minimum is finite on a
nonconstant curve.

\begin{proposition}[Curve inequality]
\label{geom:curve-inequality}
With the weights chosen above, every such curve satisfies
\begin{equation}
\label{geom:curve-bound}
 \deg_W C\ge(1+\sigma)\sum_{P\mapsto a_j}h_P,
\end{equation}
where the sum includes all branches through the prescribed points.
\end{proposition}

\begin{proof}
There is nothing to prove if the curve misses all the points. Suppose
instead that \eqref{geom:curve-bound} fails. For every sufficiently large
auxiliary integer $N$, linear algebra gives a nonzero polynomial $F_N$
of $W$-degree at most $N$ whose expansions at all the points have
$V$-order at least $(1+3\sigma)N$. Indeed, the number of its coefficients
has leading term
\[
 \frac{N^{m+1}}{(m+1)!\prod_{a=0}^m w_a},
\]
whereas the number of conditions has leading term
\[
 \frac{K(1+3\sigma)^{m+1}N^{m+1}}
      {(m+1)!\prod_{a=0}^m v_a}.
\]
Their comparison is the first inequality in \eqref{geom:sigma}. This
argument requires no independence among the conditions.

On the open set $Y\ne0$ consider the commuting frame
\begin{equation}
\label{geom:frame}
 D_0=Y\partial_Y+\sum_{i=1}^m\partial_{X_i},\qquad
 D_i=\partial_{X_i}\quad(1\le i\le m).
\end{equation}
It preserves polynomial degree at most $N$. In the formal coordinates of
\eqref{geom:jet-map},
\[
 D_0=(1+t)\partial_t,\qquad D_i=\partial_{u_i}.
\]
Thus differentiation by $D_a$ decreases weighted order by at most $v_a$.
Every $D^\gamma F_N$ with $\sum_av_a\gamma_a\le\sigma N$ retains
order at least $(1+2\sigma)N$ at each point. On a branch $P$, its order
is consequently at least $(1+2\sigma)N h_P$.

A polynomial of $W$-degree at most $N$ restricts to a rational function
on the complete model of $C$ with total pole degree at most
$N\deg_W C$. Since we assumed that \eqref{geom:curve-bound} fails, the
just obtained total zero order exceeds this pole bound. The ordinary
divisor equality for a nonzero rational function therefore forces
\begin{equation}
\label{geom:derivatives-on-curve}
 D^\gamma F_N\big|_C=0
 \qquad\left(\sum_av_a\gamma_a\le\sigma N\right).
\end{equation}
All the assertions about branch orders hold for formal power series:
each monomial has order at least its weight times $h_P$, and
cancellation can only increase the order.

We have converted excess contact into vanishing of many derivatives on
$C$. We next find a component on which that vanishing persists, as in
the derivative-ideal method of zero estimates
\cite[Section~5]{philippon1986}, and compare all its normal bases.
Put $\delta=\sigma N/L_0$. For $0\le r\le L_0$, let $V_r$ be the common
zero set in $Y\ne0$ of the derivatives of cost at most $r\delta$.
Let $C^\circ$ be the part of $C$ in this affine open set. The $V_r$
form a decreasing chain, and \eqref{geom:derivatives-on-curve} puts
$C^\circ$ in $V_{L_0}$. Choose an irreducible component of $V_{L_0}$ containing
$C^\circ$, and extend it backwards to nested irreducible components of the
earlier $V_r$. Every selected component has dimension between $1$ and
$m$: it contains a curve, and it lies in the proper zero set of
$F_N$. Since there are $L_0+1=m+3$ levels, two adjacent selected
components coincide. Denote this persistent component by $Z$, and
write $k=\operatorname{codim}Z$.

Take the derivatives of cost at most $r\delta$ as the equations of the
earlier of these two levels. Any further derivative of cost less than
$\delta$ vanishes along $Z$, because the fields commute and $Z$ is
also a component at the next level. Lemma~\ref{geom:multiplicity}, with
$\varepsilon=\sigma/L_0$, gives
\begin{equation}
\label{geom:all-bases}
 \prod_{a\in A}w_a\le C_*\prod_{b\in B}v_b
\end{equation}
for every coordinate normal basis $A$ and every frame normal basis $B$.
The equations and $Z$ may vary with $N$; the uniformity proved in
Lemma~\ref{geom:multiplicity} is exactly what permits this use.

We show that $Y$ is constant on $Z$, working on the common smooth open
subset where all the normal-basis conditions hold. Suppose first that
some positive coordinate direction is nontangent to $Z$, and choose the
largest such index $i$. Every larger positive index occurs in neither
kind of normal basis, since $D_j=\partial_{X_j}$ for $j>0$.
Some coordinate normal basis $A$ contains $i$. If a frame normal basis
$B$ omitted $i$, the largest positive index where these sets differ
would be $i$. Then
\eqref{geom:separation} would contradict \eqref{geom:all-bases}.
Thus every frame normal basis contains $D_i$.

The other frame vectors consequently fail to span the normal quotient.
Their ambient span is the hyperplane
\[
 \ker(dX_i-dY/Y).
\]
For an ambient hyperplane $Q$ and a tangent space $T$, the image of
$Q$ in the quotient by $T$ fails to be surjective precisely when
$T\subseteq Q$. We obtain
\begin{equation}
\label{geom:log-differential}
 dX_i=dY/Y\quad\text{on }Z.
\end{equation}
This identity forces $Y$ to be constant on the algebraic variety $Z$.
Indeed, if $Y$ were nonconstant, choose a smooth point where $dY\ne0$
and intersect with general algebraic hyperplanes through that point
until a curve remains on which $Y$ is nonconstant. On the smooth
complete model of this curve, \eqref{geom:log-differential} is an
identity of meromorphic differentials. A nonconstant meromorphic
function $Y$ has a zero or pole, where $dY/Y$ has nonzero residue
$\operatorname{ord}_P Y$. An exact differential $dX_i$ has zero
residue at every point. This is a contradiction. It follows also
that $X_i$ is constant, since $dX_i=0$ in characteristic zero.

If instead every positive coordinate direction is tangent, the $m$
independent vectors $\partial_{X_i}$ all lie in the tangent space of
$Z$. Properness forces that tangent space to be their span, and
$dY=0$ on $Z$ again. We conclude in all cases that $Y$ is constant
on $Z$, and hence on $C$. Since the nonzero numbers $y_j$ are pairwise
distinct, $C$ passes through at most one prescribed center, say $a_j$.
At least one $X_l$ is nonconstant on $C$, since $C$ is a curve and $Y$
is constant. On every branch through this center, $Y=y_j$ and therefore
$\log(Y/y_j)=0$. The nonzero rational function $X_l-c_{jl}$ has total
pole degree at most $w_l\deg_W C$, whereas its order at each of these
branches is at least $(w_l/\theta)h_P$. Consequently
\[
 \deg_W C\ge\theta^{-1}\sum_{P\mapsto a_j}h_P
           >(1+\sigma)\sum_{P\mapsto a_j}h_P.
\]
This contradicts the assumed failure of the curve inequality.
The original proof's second argument inside the fiber $Y=1$ is
unnecessary here. The hypothesis $K\theta^m<1$ has been retained to
keep the subsequent parameter choice identical to the source proof.
\end{proof}

\subsection{From the curve inequality to all coefficient packets}

We complete the proof of Theorem~\ref{geom:interpolation}. The argument uses
ordinary projective schemes and the ordinary blow-up; smoothness or
normality of the compactification will not be needed. Curve positivity
and asymptotic jet generation are closely related in the Seshadri-constant
approach of \cite[Section~6]{demailly1992}. Here we supply the passage
for our weighted ideals at several points, including the eventual
ordinary-power comparison required on a possibly singular scheme.

Choose an integer $R$ sufficiently large that $R/w_a$ and $R/v_a$ are
positive integers for every index. Embed affine space by all monomials
of $W$-degree at most $R$, including the constant monomial, and let
$X$ be the projective closure of the image. Since these monomials
include each coordinate, the chart of the constant monomial is
precisely the original affine space. Let $L=\mathcal O_X(1)$, a very
ample line bundle.

For any curve meeting this chart,
\begin{equation}
\label{geom:projective-degree}
 \deg(L|_C)=R\deg_W C.
\end{equation}
To see this on its smooth complete model, at each point clear the
largest pole of the homogeneous monomial coordinates. Every monomial
has pole order at most $R$ times the maximum in
\eqref{geom:curve-degree}. A pure power of one coordinate attains
that bound whenever the maximum is positive, and the constant
monomial attains it when the maximum is zero. The resulting sections
have no common zero, proving \eqref{geom:projective-degree}.

To make the local data algebraic, replace $\log(1+t)$ in each
coordinate by a fixed Taylor polynomial $G_i(t)$ whose omitted
terms have $t$-weight strictly greater than $v_i$. Then
\[
 t=Y/y_j-1,\qquad u_i'=X_i-c_{ji}-G_i(t)
\]
are regular local coordinates at $a_j$. This replacement preserves
$h_P$. Indeed, if $t$ is not identically zero on a branch, each
omitted term has order divided by $v_i$ strictly larger than
$\operatorname{ord}_P(t)/v_0$; if $t$ is identically zero, its tail
vanishes. It also preserves interpolation surjectivity for every
threshold: the formal coordinate change fixes $t$ and changes
$u_i$ only by terms of weight greater than $v_i$. It and its inverse
preserve the weighted filtration and induce inverse triangular maps
on each finite coefficient packet.

At $a_j$ form the finite-colength ideal
\[
 I_j=\left(t^{R/v_0},(u_1')^{R/v_1},\ldots,(u_m')^{R/v_m}\right).
\]
The points are distinct. These local ideals, together with the unit
ideal away from the points, define a coherent ideal sheaf $I$ on $X$.
For a branch through a center,
\begin{equation}
\label{geom:ideal-valuation}
 \operatorname{ord}_P I
 =\min_a\left\{(R/v_a)\operatorname{ord}_P z_a\right\}
 =Rh_P,
 \qquad z_0=t,\quad z_i=u_i'.
\end{equation}
Let $p:X'\to X$ be the ordinary blow-up of $I$, namely the relative
Proj of its Rees algebra, and write
\[
 I\mathcal O_{X'}=\mathcal O_{X'}(-E).
\]
The scheme $X$ is integral and Noetherian, $I$ is a coherent nonzero
ideal, and $E$ is an effective Cartier divisor. The tautological
bundle $\mathcal O_{X'}(-E)$ is relatively ample; these are the
standard blow-up properties in
\cite[Tags~02OS, 02ND, and~02NS]{stacks}.

Write $A=p^*L$ in additive divisor notation. We claim that
\begin{equation}
\label{geom:nef-class}
 A-(1+\sigma)E\quad\text{is nef}.
\end{equation}
For a noncontracted integral curve $\Gamma$ upstairs whose image
meets the affine chart, the map to its image is birational, since
$p$ is an isomorphism away from the finite center. On the complete
normalizations, \eqref{geom:projective-degree} and
\eqref{geom:ideal-valuation} give
\[
 A\cdot\Gamma=R\deg_W p(\Gamma),\qquad
 E\cdot\Gamma=R\sum_{P\mapsto a_j}h_P.
\]
Proposition~\ref{geom:curve-inequality} gives the desired nonnegative
degree. A curve whose image misses the affine chart avoids the
exceptional divisor. A contracted curve has $A$-degree zero and
positive degree for $-E$ by relative ampleness. These cases account
for every integral curve, proving \eqref{geom:nef-class}.

For sufficiently large $a>1$, the divisor $aA-E$ is ample. One
can see this directly: choose $b$ so that $I\otimes L^b$ is globally
generated. The Rees construction embeds $X'$ in a product
$X\times\mathbb P^r$, with the second hyperplane bundle restricting
to $bA-E$. Tensoring with the first very ample bundle makes
$(b+1)A-E$ ample. Now the identity
\begin{equation}
\label{geom:ample-combination}
 A-E=
 \frac{a-1}{a(1+\sigma)-1}\bigl(A-(1+\sigma)E\bigr)
 +\frac{\sigma}{a(1+\sigma)-1}(aA-E)
\end{equation}
expresses $A-E$ as a positive multiple of a nef class plus a
positive multiple of an ample class. It is therefore ample;
see \cite[Corollary~1.4.10]{lazarsfeld}. This theorem applies to
projective schemes and does not require a smooth blow-up.

For all sufficiently large $n$, the graded-piece and relative
vanishing theorems for the ordinary Rees algebra give
\begin{equation}
\label{geom:rees-pushdown}
 p_*\mathcal O_{X'}(-nE)=I^n,\qquad
 R^q p_*\mathcal O_{X'}(-nE)=0\quad(q>0).
\end{equation}
Indeed, on any affine open of the Noetherian base the Rees algebra
is a Noetherian graded algebra generated in degree one. Its
sufficiently high graded pieces recover the sections of its
associated sheaf on Proj, and the higher cohomology of those twists
vanishes. A finite affine cover gives one common threshold;
see \cite[Tags~0AG6 and~0AG7]{stacks}. In particular,
\eqref{geom:rees-pushdown} uses ordinary powers $I^n$, not their
integral closures, and makes no assertion about small $n$.

By the projection formula, Leray, and Serre vanishing for the ample
line bundle $A-E$,
\begin{equation}
\label{geom:serre-vanishing}
 H^1(X,I^n\otimes L^n)
 =H^1(X',\mathcal O_{X'}(n(A-E)))=0
\end{equation}
for all sufficiently large $n$; see \cite[Tag~0B5U]{stacks}.
The exact sequence for $I^n$ consequently gives a surjection from
$H^0(X,L^n)$ onto all its local coefficient data modulo $I^n$.
The constant projective coordinate trivializes $L$ on the affine
chart. Finite-colength local quotients are unchanged by completion,
so these algebraic classes are also the formal coefficient data used
in the theorem. In these trivializations, every element of $I^n$ has
weighted order at least $nR$, because every displayed generator
of $I$ has weight $R$. Thus
\begin{equation}
\label{geom:ideal-inclusion}
 I^n\subseteq\{\text{local series of weight at least }nR\}.
\end{equation}
Surjectivity modulo $I^n$ therefore implies surjectivity onto the
smaller quotient consisting of the desired packets of weight less
than $nR$. Equality in \eqref{geom:ideal-inclusion} is not needed.

Finally, every section of $L^n$ for sufficiently large $n$ is
supplied by a homogeneous degree-$n$ polynomial in the projective
embedding coordinates. This follows from Serre vanishing applied
to the homogeneous ideal sheaf of $X$ in its projective space.
Such a polynomial restricts on the affine chart to one of
$W$-degree at most $nR$. Taking $H=nR$ and undoing the filtered
Taylor-polynomial coordinate change proves the surjectivity in
\eqref{geom:jet-map}, and completes the proof of
Theorem~\ref{geom:interpolation}.

\section{The arithmetic and analytic determinant estimates}
\label{det:section}

Interpolation now supplies the algebraic rank needed for an
interpolation determinant, in the sense of
\cite[Section~6]{laurent1991}. We estimate one full-row minor from
below arithmetically, and then translate its rows to exact logarithmic
points to obtain an analytic upper bound.

Throughout this section all parameters except $H$ are fixed.  Write
\[
 w\alpha=\sum_{i=1}^m w_i\alpha_i,
 \qquad w_* =\min_{1\le i\le m}w_i,
 \qquad \omega=\log 2.
\]
Let $\nu>2$, and suppose that integers $p_i$ and $q_i\ge2$ satisfy
\begin{equation}
 \left|\log 2-\frac{p_i}{q_i}\right|\le q_i^{-\nu},
 \qquad w_i=\lceil\log q_i\rceil,
 \qquad r_i=\frac{p_i}{q_i}
 \quad(1\le i\le m).
 \label{det:approximations}
\end{equation}
Assume that $p_i\ne0$ and that $m,K,w_0,v_0,\theta,w_1,\ldots,w_m$
satisfy the hypotheses of Theorem~\ref{geom:interpolation}, with centers
$y_j=2^j$ and $c_{ji}=jr_i$, $0\le j<K$.  In particular $w_i$ is a positive integer
for $i>0$, while $w_0$ and $v_0$ are positive rational numbers.  Fix
$F_0>2/\theta$, and put
\begin{equation}
 T_i=\left\lceil\frac{F_0w_i}{v_0}\right\rceil,
 \qquad
 G_i(t)=\sum_{1\le k<T_i}\frac{(-1)^{k+1}t^k}{k}.
 \label{det:truncations}
\end{equation}
Empty sums are zero.  The first omitted term has weight
$v_0T_i\ge F_0w_i>w_i/\theta$.  Consequently the substitution
$u_i\mapsto u_i+G_i(t)-\log(1+t)$ is an invertible substitution preserving
the weighted filtration.  It induces an invertible map on every finite
space of coefficients of weight less than $H$.  Thus replacing the
logarithms in Theorem~\ref{geom:interpolation} by these $G_i$ preserves
surjectivity.

\subsection{A full-row-rank minor and its arithmetic size}

Use the monomials
\[
 P(Y,X)=Y^hX^\alpha,
 \qquad h\in\ZZ_{\ge0},\quad \alpha\in\ZZ_{\ge0}^m,
 \qquad w_0h+w\alpha\le H
\]
as columns, and the triples
\[
 \rho=(j,s,\beta),\qquad
 0\le j<K,\quad s\in\ZZ_{\ge0},\quad\beta\in\ZZ_{\ge0}^m,
 \qquad v_0s+\frac{w\beta}{\theta}<H
\]
as rows.  The matrix entry is the coefficient of $t^su^\beta$ in
\begin{equation}
 P\bigl(2^j(1+t),jr_1+G_1(t)+u_1,\ldots,jr_m+G_m(t)+u_m\bigr).
 \label{det:matrix}
\end{equation}
For all sufficiently large, sufficiently divisible $H$, choose a square
minor $\Delta_H\ne0$ using every row, as permitted by
Theorem~\ref{geom:interpolation}.  Denote its size by $M=M_H$, and set
\begin{equation}
 \bar b=\bar b_H=\frac{1}{MH}\sum_{\rho=(j,s,\beta)}w\beta.
 \label{det:bbar}
\end{equation}
The elementary lattice count for a simplex gives
\begin{equation}
 M\sim
 \frac{K\theta^mH^{m+1}}{(m+1)!\,v_0\prod_{i=1}^m w_i},
 \qquad 0\le\bar b\le\theta.
 \label{det:row-count}
\end{equation}
Indeed, after scaling by $H$, the lattice boxes of side lengths $H^{-1}$
are Riemann sums for the fixed simplex.  Its volume is
$\theta^m/((m+1)!v_0\prod_iw_i)$; its boundary has volume zero, so the
strict inequality in the row definition does not affect the leading
term.  We always let $H$ tend to infinity only after all the weights and
centers have been fixed.

\begin{lemma}[Arithmetic lower bound]\label{det:arithmetic}
Let $\Lambda=4\log2$, and define
\begin{equation}
 E_{\rm ar}
 =\frac{\Lambda F_0m}{v_0}
  +\Lambda\sum_{i=1}^m\frac1{w_i}+\frac{\theta}{w_*}.
 \label{det:arith-error}
\end{equation}
Then every minor just chosen satisfies
\begin{equation}
 \frac{\log|\Delta_H|}{MH}
 \ge -(1-\bar b)-E_{\rm ar}.
 \label{det:lower}
\end{equation}
\end{lemma}

\begin{proof}
Let $L_i=\operatorname{lcm}(1,\ldots,T_i-1)$, interpreted as $1$ if
$T_i=1$, and put
\[
 D_H=\prod_{i=1}^m L_i^{\lfloor H/w_i\rfloor}.
\]
Scale a column $Y^hX^\alpha$ by $\prod_iq_i^{\alpha_i}$ and a row
$(j,s,\beta)$ by $\prod_iq_i^{-\beta_i}$.  The scaled entry is zero
unless $\alpha\ge\beta$ componentwise; otherwise it is
\[
 \prod_i\binom{\alpha_i}{\beta_i}
 2^{jh}[t^s](1+t)^h
 \prod_i\bigl(q_i(jr_i+G_i(t))\bigr)^{\alpha_i-\beta_i}.
\]
Here $q_ijr_i=jp_i\in\ZZ$.  Since $G_i$ has coefficients in
$L_i^{-1}\ZZ$, the denominator of this entry divides
$\prod_iL_i^{\alpha_i-\beta_i}$, which divides $D_H$.  Thus multiplying
every entry of the scaled matrix by $D_H$ gives a matrix over
$\ZZ$ with nonzero determinant.  A nonzero integer has
modulus at least one, so
\begin{equation}
 \begin{aligned}
 \log|\Delta_H|
 &\ge -M\log D_H
      -\sum_{\text{columns }\alpha}\sum_i\alpha_i\log q_i\\
 &\quad+\sum_{\text{rows }\beta}\sum_i\beta_i\log q_i.
 \end{aligned}
 \label{det:clearing}
\end{equation}
For each selected column, $\sum_i\alpha_i\log q_i\le w\alpha\le H$.
Also $0\le w_i-\log q_i<1$, and hence
\[
 \sum_{\text{rows }\beta}\sum_i\beta_i\log q_i
 \ge MH\bar b-\sum_{\text{rows }\beta}\sum_i\beta_i
 \ge MH\bar b-\frac{MH\theta}{w_*}.
\]

For completeness, $\log\operatorname{lcm}(1,\ldots,n)\le\Lambda n$.
To see this, the factorial valuation formula shows that each prime
power in $(\ell,2\ell]$ contributes its logarithmic prime weight to
$\log\binom{2\ell}{\ell}$; all the other contributions to the latter
are nonnegative.  Thus the increase of the log least common multiple
across this interval is at most $2\ell\log2$.  Summing over dyadic
intervals up to the smallest power of two above $n$ gives at most
$4n\log2$.  Therefore
\[
 \frac{\log D_H}{H}
 \le\Lambda\sum_i\frac{T_i}{w_i}
 \le\frac{\Lambda F_0m}{v_0}+\Lambda\sum_i\frac1{w_i}.
\]
Substitution in \eqref{det:clearing} proves \eqref{det:lower}.
\end{proof}

\subsection{Translation to entire functions}

For every multi-index $a\in\ZZ_{\ge0}^m$ and every selected column $P$,
define the entire function
\begin{equation}
 f_{a,P}(z)=[u^a]P(e^z,z+u_1,\ldots,z+u_m).
 \label{det:test-functions}
\end{equation}
We regard $(f_{a,P})_P$ as a row vector of entire functions, indexed by
the selected columns.  For $P=Y^hX^\alpha$ it has the explicit entry
\begin{equation}
 f_{a,P}(z)=
 \begin{cases}
 \displaystyle\binom{\alpha}{a}e^{hz}z^{|\alpha|-|a|},&a\le\alpha,\\
 0,&\text{otherwise}.
 \end{cases}
 \label{det:test-monomials}
\end{equation}
In particular only indices $a$ with $wa\le H$ can contribute.

\begin{lemma}[Exact row translation]\label{det:translation}
Every row $\rho=(j,s,\beta)$ of \eqref{det:matrix} is a sum of at most
\begin{equation}
 Q_H=(\lfloor H\rfloor+1)^{2m}
          (\lfloor H/v_0\rfloor+1)
 \label{det:term-count}
\end{equation}
scalar multiples of row vectors of the form
\[
 \bigl([t^\ell]f_{a,P}(j\omega+\log(1+t))\bigr)_P,
 \qquad a\ge\beta,\quad wa\le H,\quad0\le\ell\le s.
\]
Each scalar $\xi$ occurring with index $a$ satisfies
\begin{equation}
 \begin{aligned}
 |\xi|&\le\exp\{-\nu w(a-\beta)+H E_{\rm tr}\},\\
 E_{\rm tr}
 &=\frac{\nu}{F_0}+\frac{\log2}{v_0}
   +\frac{\log4+\log(2K)+\nu}{w_*}.
 \end{aligned}
 \label{det:translation-bound}
\end{equation}
\end{lemma}

\begin{proof}
Put
\[
 \epsilon_i=j(r_i-\omega),\qquad
 \tau_i(t)=G_i(t)-\log(1+t).
\]
Since $e^{j\omega}=2^j$, the substitutions in \eqref{det:matrix} can be
written as $Y=e^z$, $X_i=z+u_i+\epsilon_i+\tau_i(t)$, where
$z=j\omega+\log(1+t)$.  Expanding first in the variables $u_i$, then
in the constants $\epsilon_i$ and tails $\tau_i$, gives the exact row
vector identity
\begin{align}
 \text{row}_{j,s,\beta}
 ={}&\sum_{\substack{a\ge\beta\\wa\le H}}
       \binom a\beta
       \sum_{0\le d\le a-\beta}\binom{a-\beta}{d}
       \epsilon^{a-\beta-d}
       \sum_{k=0}^s
       \left([t^k]\prod_i\tau_i(t)^{d_i}\right)
 \nonumber\\[-2pt]
 &\hspace{35mm}\cdot
 \bigl([t^{s-k}]f_{a,P}(j\omega+\log(1+t))\bigr)_P.
 \label{det:row-identity}
\end{align}
All scalar factors in this identity are independent of the column
$P$.  This is therefore an identity to which determinant
multilinearity applies directly.

The product of tails vanishes to order at least $\sum_i d_iT_i$.
Consequently a nonzero coefficient in \eqref{det:row-identity}
satisfies
\begin{equation}
 \sum_i d_iw_i\le\frac{v_0}{F_0}\sum_i d_iT_i
 \le\frac{v_0s}{F_0}<\frac{H}{F_0}.
 \label{det:tail-budget}
\end{equation}
On $|t|=1/2$,
\[
 |\tau_i(t)|
 \le\sum_{k\ge T_i}\frac{2^{-k}}k\le1.
\]
Cauchy's coefficient estimate thus bounds the tail coefficient in
\eqref{det:row-identity} by $2^k\le2^s$.  From
\eqref{det:approximations},
\[
 |\epsilon_i|\le2Kq_i^{-\nu}
 \le2K e^\nu e^{-\nu w_i}.
\]
Moreover
$\binom a\beta\binom{a-\beta}{d}\le4^{|a|}$ and
$|a|\le H/w_*$.  Combining these inequalities and using
\eqref{det:tail-budget} gives
\[
 \log|\xi|
 \le-\nu w(a-\beta)+\frac{\nu H}{F_0}
       +\frac{H\log2}{v_0}
       +\frac{H(\log4+\log(2K)+\nu)}{w_*}
\]
for each nonzero scalar.  A zero scalar obeys the required bound as
well.  Finally, each coordinate of $a$ and $d$ is at most $H$ because
$w_i\ge1$, while $0\le k\le H/v_0$.  This proves the term count and
the Lemma.
\end{proof}

\subsection{Repeated Taylor degrees force a quadratic saving}

The Taylor expansion of interpolation determinants eliminates repeated
degrees \cite[Theorem~1]{laurent2000}. Our rows fall into different
transverse-index groups, so this cancellation is used only inside a
fixed group. The next bound is uniform over all row choices.

\begin{lemma}[Collision estimate]\label{det:collision}
Choose one translated test row from each row expansion in
Lemma~\ref{det:translation}, and remove its scalar factor.  Let $T$ be
the resulting square matrix, and let $n_a$ be the number of its rows
with index $a$.  With
\begin{equation}
 c=\frac{\log2}{4},\qquad
 E_{\rm hol}=\frac{100K}{w_0}+\frac{\log2}{v_0}
             +\frac{\log(200K)}{w_*},
 \label{det:hol-error}
\end{equation}
there is a quantity $\rho_H\to0$, independent of all these row
choices and of the selected columns, such that
\begin{equation}
 |\det T|
 \le\exp\left\{-c\sum_a n_a^2+MH(E_{\rm hol}+\rho_H)\right\}.
 \label{det:collision-bound}
\end{equation}
\end{lemma}

\begin{proof}
Put $R=100K$.  Equation~\eqref{det:test-monomials} gives, on $|z|\le R$,
\[
 |f_{a,P}(z)|\le
 \exp\left\{H\left(\frac R{w_0}
                      +\frac{\log(2R)}{w_*}\right)\right\}
 =:\mathcal D_H.
\]
Expand
\[
 f_{a,P}(z)=\sum_{d\ge0}c_{a,d,P}(z/R)^d.
\]
Cauchy's estimate gives $|c_{a,d,P}|\le\mathcal D_H$.  A row of $T$ is
a test of order $\ell\le s$ at some center $j$.  Since
\[
 |j\omega+\log(1+t)|
 \le(K-1)\log2+\log2<50K=R/2
 \qquad(|t|=1/2),
\]
its value on $(z/R)^d$ has modulus at most $2^\ell2^{-d}$.

Expand the determinant in these power-series rows.  The expansion is
absolutely convergent: a coefficient determinant has modulus at most
$M!\mathcal D_H^M$, and the sum of the product bounds $2^{-d}$ over
all degrees is finite.  If two rows have the same pair $(a,d)$, their
coefficient vectors $(c_{a,d,P})_P$ agree, and that summand vanishes.
Thus every nonzero summand has distinct degrees within each group
$a$, and so
\[
 \sum_{\text{rows}}d\ge\sum_a\binom{n_a}{2}.
\]
Use half of the geometric decay for this lower bound and sum the
other half freely over all nonnegative degrees.  Since
$\sum\ell\le MH/v_0$, this gives
\[
 |\det T|
 \le M!\mathcal D_H^M2^{MH/v_0}
       2^{-\frac12\sum_a\binom{n_a}{2}}
       (1-2^{-1/2})^{-M}.
\]
Because $\sum_a n_a=M$, the asserted inequality follows with, for
example,
\[
 \rho_H=\frac{\log M+\frac14\log2
                    -\log(1-2^{-1/2})}{H}.
\]
Equation~\eqref{det:row-count} shows that $\rho_H\to0$ with the fixed
parameters stipulated above.
\end{proof}

\subsection{A comparison with two alternatives}

Put
\begin{equation}
 E_{\rm an}=E_{\rm tr}+E_{\rm hol}.
 \label{det:analytic-error}
\end{equation}

Let $0<A<1$ set a cutoff $AH$ for transverse weights, and let
$0<\eta<1$ be a fraction of rows. The main arithmetic cost is
$1-\bar b$. In each term of the determinant expansion, if at least
$\eta M$ rows have $wa\le AH$, they give a collision saving;
otherwise the rows with $wa>AH$ supply approximation-error factors.
We seek a saving larger than $1-\bar b+E_{\rm ar}+E_{\rm an}$, and
choose $A$ and $\eta$ in Section~\ref{sec:parameters}.

\begin{proposition}[Determinant comparison]\label{det:comparison}
Let $0<A<1$ and $0<\eta<1$, and define
\begin{equation}
 g=\nu\bigl(A(1-\eta)-\theta\bigr)-(1-\theta),
 \qquad
 L=\frac{\eta^2K\theta^m}{(m+1)v_0A^m}.
 \label{det:budgets}
\end{equation}
The approximations \eqref{det:approximations} and the interpolation
hypotheses cannot hold simultaneously if
\begin{equation}
 g>0,\qquad
 E_{\rm ar}+E_{\rm an}<g,
 \qquad
 cL>1+E_{\rm ar}+E_{\rm an}.
 \label{det:contradiction-conditions}
\end{equation}
\end{proposition}

\begin{proof}
Let
\[
 D_A(H)=\#\{a\in\ZZ_{\ge0}^m:wa\le AH\},
 \qquad L_H=\frac{\eta^2M}{H D_A(H)}.
\]
The simplex count, now in dimension $m$, and
\eqref{det:row-count} show that
\[
 D_A(H)\sim\frac{(AH)^m}{m!\prod_iw_i},
 \qquad L_H\longrightarrow L.
\]
Expand $\Delta_H$ by the row identities in
Lemma~\ref{det:translation}.  For a choice of terms with indices
$a_\rho$, the product of their scalar factors has modulus at most
\[
 \exp\left\{-\nu\sum_\rho wa_\rho
             +\nu MH\bar b+MH E_{\rm tr}\right\}.
\]
There are at most $Q_H^M$ term choices.  Apply
Lemma~\ref{det:collision} to each one.

Suppose first that at least $\eta M$ of its indices obey $wa\le AH$.
There are at most $D_A(H)$ such indices, so Cauchy--Schwarz gives
\[
 \sum_a n_a^2\ge\frac{(\eta M)^2}{D_A(H)}=MH L_H.
\]
The scalar main exponent $-\nu\sum_\rho w(a_\rho-\beta_\rho)$ is
nonpositive, because $a_\rho\ge\beta_\rho$.  Discarding it leaves
an upper bound $-cL_H+E_{\rm an}+\rho_H$ after division by $MH$.

In the other case, more than $(1-\eta)M$ indices have $wa>AH$.
Hence $\sum_\rho wa_\rho\ge MH A(1-\eta)$.  Discarding the
nonpositive collision term leaves the upper bound
$-\nu(A(1-\eta)-\bar b)+E_{\rm an}+\rho_H$.
Taking absolute values before summing the determinant expansion gives
\begin{equation}
 \begin{aligned}
 \frac{\log|\Delta_H|}{MH}
 &\le E_{\rm an}+\varepsilon_H+
       \max\{-cL_H,-\nu(A(1-\eta)-\bar b)\},\\
 \varepsilon_H&=\rho_H+\frac{\log Q_H}{H}\longrightarrow0.
 \end{aligned}
 \label{det:upper}
\end{equation}
All bounds here are uniform over the row-term choices.

For the second alternative, the gap between the two main exponents
is at least $g$, since $\bar b\le\theta$ gives
\[
 \nu(A(1-\eta)-\bar b)-(1-\bar b)
 =\nu A(1-\eta)-1-(\nu-1)\bar b\ge g.
\]
Under \eqref{det:contradiction-conditions}, both quantities inside the
maximum in \eqref{det:upper}, after adding
$E_{\rm an}+\varepsilon_H$, are strictly less than
$-(1-\bar b)-E_{\rm ar}$ for all sufficiently large $H$.  For the
first quantity, use $L_H\to L$ and $1-\bar b\le1$; for the second,
use the displayed uniform gap.  This contradicts
Lemma~\ref{det:arithmetic}.
\end{proof}

\section{Choice of parameters and proof of Theorem~\ref{thm:main}}
\label{sec:parameters}

We now combine the interpolation theorem and the determinant estimates.
The order of choices is essential: the dimension is fixed before the
approximation scales, and all these data are fixed before the polynomial
degree tends to infinity.

\begin{lemma}\label{par:constants}
For every real \(\nu>2\), there are positive rational numbers
\(\theta,A,B,C\) such that
\begin{equation}\label{par:inequalities}
0<\theta<A<B<1,\qquad
\nu(A-\theta)>1-\theta,\qquad
C>1,\quad B<1/C,\quad B<C\theta<1.
\end{equation}
One can then choose \(0<\eta<1\) such that
\begin{equation}\label{par:gap}
g:=\nu\bigl(A(1-\eta)-\theta\bigr)-(1-\theta)>0.
\end{equation}
\end{lemma}

\begin{proof}
Choose a rational \(b\) with \(1/2<b<1-1/\nu\).
For a sufficiently small positive rational \(\delta\), put
\(\theta=1-\delta\) and \(A=1-b\delta\). Then \(0<\theta<A<1\), and
\[
\nu(A-\theta)-(1-\theta)
=\bigl(\nu(1-b)-1\bigr)\delta>0,
\]
whereas
\[
\theta-A^2=(2b-1)\delta-b^2\delta^2>0.
\]
It follows that the interval \((A/\theta,1/A)\) is nonempty.
Choose a rational \(C\) in this interval. Since \(A>\theta\), we have
\(C>1\); since \(\theta<A\), we also have
\(C\theta<\theta/A<1\).
Now choose a rational \(B\) between \(A\) and
\(\min(1/C,C\theta)\). This proves \eqref{par:inequalities}.
Finally, the strict inequality at \(\eta=0\) allows a sufficiently
small \(\eta>0\) in \eqref{par:gap}.
\end{proof}

The two inequalities used in this construction have different roles.
The inequality \(\nu(A-\theta)>1-\theta\) supplies the
approximation-error gap. The inequality \(A^2<\theta\) permits
\(A/\theta<C<1/A\), and hence a choice of \(B\) for which
\(CB<1\), \(C\theta/B>1\), and \(B/A>1\). With the full-simplex
counts of Section~\ref{det:section}, these three ratios make the
dimension-dependent errors tend to zero while the collision saving
tends to infinity. We now choose one finite dimension realizing
both requirements; it will remain fixed throughout the degree limit.

\begin{proof}[Proof of Theorem~\ref{thm:main}]
Fix \(\nu>2\). Suppose for a contradiction that there are arbitrarily
large positive integers \(q\) for which some \(p\in\ZZ\) satisfies
\begin{equation}\label{par:good}
\left|\log 2-\frac pq\right|\le q^{-\nu}.
\end{equation}
Fix \(\theta,A,B,C,\eta,g\) from Lemma~\ref{par:constants}, and put
\(\varepsilon_0=\min(g/2,1/2)\).
First choose \(F_0>2/\theta\) so large that
\(\nu/F_0<\varepsilon_0/3\).

For an integer \(m\), set
\begin{equation}\label{par:dimension}
K=\lfloor C^m\rfloor,\qquad
w_0=B^{-m},\qquad
v_0=2K\theta^m w_0.
\end{equation}
These are admissible rational interpolation data:
\begin{equation}\label{par:volume}
K\theta^m\le(C\theta)^m<1,\qquad
K\frac{w_0}{v_0}\theta^m=\frac12.
\end{equation}
The relevant limits as \(m\to\infty\) are
\begin{equation}\label{par:limits}
\frac{K}{w_0}\le(CB)^m\longrightarrow0,\qquad
v_0\sim2\left(\frac{C\theta}{B}\right)^m\longrightarrow\infty,
\qquad
\frac{m}{v_0}\longrightarrow0.
\end{equation}
With the notation of Proposition~\ref{det:comparison}, these choices give
\begin{equation}\label{par:collision-growth}
L=\frac{\eta^2K\theta^m}{(m+1)v_0A^m}
 =\frac{\eta^2(B/A)^m}{2(m+1)}\longrightarrow\infty.
\end{equation}
Thus, with \(\Lambda=4\log2\) and \(c=(\log2)/4\), choose \(m\) large
enough that
\begin{equation}\label{par:dimension-margin}
\frac{\Lambda F_0m+2\log2}{v_0}+\frac{100K}{w_0}
 <\frac{\varepsilon_0}{3},
\qquad cL>2.
\end{equation}
Fix this \(m\), and hence \(K,w_0,v_0\).

The remaining constants are now fixed. Choose approximations
\(p_i/q_i\) satisfying \eqref{par:good}, with
\(w_i=\lceil\log q_i\rceil\), successively sufficiently large for
Theorem~\ref{geom:interpolation} and, in addition, for
\begin{equation}\label{par:weight-margin}
\Lambda\sum_{i=1}^m\frac1{w_i}
+\frac{\theta+\log4+\log(2K)+\nu+\log(200K)}{w_*}
 <\frac{\varepsilon_0}{3},
\qquad w_*=\min_{i>0}w_i.
\end{equation}
This is possible: first impose a sufficiently large lower bound on
every \(w_i\), using \(\sum_i1/w_i\le m/w_*\), and then impose the
successive geometric thresholds. The assumed denominators are
unbounded, so every finite sequence of such requirements can be met.
The thresholds for weight separation are independent of the centers.
Taking each \(q_i\) sufficiently large ensures \(p_i\ne0\), since
\(q_i^{-\nu}<\log2\). The centers are \(y_j=2^j\), \(c_{ji}=jr_i\), and \(r_i=p_i/q_i\).
The \(y_j\) are nonzero and pairwise distinct, as required by the
modified interpolation theorem. The thresholds are independent of
both the \(y_j\) and \(c_{ji}\).

Using \eqref{det:arith-error}, \eqref{det:translation-bound},
\eqref{det:hol-error}, and \eqref{det:analytic-error}, the sum of errors
is exactly
\[
\begin{aligned}
E_{\rm ar}+E_{\rm an}
&=\frac{\nu}{F_0}
 +\frac{\Lambda F_0m+2\log2}{v_0}+\frac{100K}{w_0}\\
&\quad +\Lambda\sum_i\frac1{w_i}
 +\frac{\theta+\log4+\log(2K)+\nu+\log(200K)}{w_*}.
\end{aligned}
\]
Our three choices make this less than \(\varepsilon_0\).
In particular it is less than \(g\), while
\(cL>2>1+E_{\rm ar}+E_{\rm an}\).
All hypotheses of Proposition~\ref{det:comparison} now hold,
a contradiction. The degree \(H\) in that Proposition tends to infinity
only after the dimension, weights, and centers just chosen are fixed.

We have excluded \eqref{par:good} for arbitrarily large \(q\), for
each fixed \(\nu>2\). This gives the required eventual lower bound,
in fact with a strict inequality. The contradiction allowed unreduced
fractions and zero error. If \(\log 2\) were rational, exact representations
with arbitrarily large denominators would contradict it; hence
\(\log 2\) is irrational.

Finally divide the unit interval into \(N\) equal parts and apply
pigeonhole to \(0,\log 2,\ldots,N\log 2\) modulo one. There exist \(p\in\ZZ\)
and \(1\le q\le N\) such that
\[
0<\left|\log 2-\frac pq\right|\le\frac1{qN}\le\frac1{q^2}.
\]
After reduction, the same error is at most the inverse square of the
reduced denominator. These errors tend to zero as \(N\to\infty\);
irrationality forces the reduced denominators to be unbounded.
Every exponent less than two therefore occurs infinitely often in
the defining inequality. Hence \(\mu(\log 2)\ge2\), and the upper bound
already proved gives \(\mu(\log 2)=2\).
\end{proof}



% SPDX-License-Identifier: Apache-2.0
% Candidate extension; insert after the log 2 parameter argument.
\section{Extension to positive rational logarithms}
\label{rat:section}

This section belongs to the candidate derivation.  Its principal
conclusion and its corollaries are conditional on
Theorem~\ref{geom:interpolation}, including the assertion that its
successive weight thresholds are independent of the interpolation
centers.  The following argument records all changes needed for a
positive rational base; it does not certify the proposed geometric
interpolation theorem or claim a new unconditional result.

\begin{theorem}[Conditional rational-logarithm extension; candidate]
\label{rat:main}
Assume Theorem~\ref{geom:interpolation} with its stated quantifiers.
For every fixed $a\in\QQ_{>0}\setminus\{1\}$ and every real $\nu>2$,
there is an integer $Q(a,\nu)$ such that
\begin{equation}
 \left|\log a-\frac pq\right|>q^{-\nu}
 \qquad(p\in\ZZ,\ q\in\ZZ,\ q\ge Q(a,\nu)).
 \label{rat:eventual-bound}
\end{equation}
Here fractions need not be reduced, and $\log$ is the real natural
logarithm.  In particular, conditionally on that interpolation
statement, $\log a$ is irrational and $\mu(\log a)=2$.
\end{theorem}

We first give the determinant estimates for this fixed base.  Write
$a=u/v$ in lowest terms, with positive integers $u,v$, and put
$\omega=\log a$.  Fix interpolation data
$m,K,w_0,v_0,\theta,w_1,\ldots,w_m$ satisfying the hypotheses of
Theorem~\ref{geom:interpolation}.  Suppose that
\begin{equation}
 \left|\omega-\frac{p_i}{q_i}\right|\le q_i^{-\nu},\qquad
 q_i\ge2,\qquad w_i=\lceil\log q_i\rceil,\qquad
 r_i=p_i/q_i \quad(1\le i\le m).
 \label{rat:approximations}
\end{equation}
Set $w\alpha=\sum_iw_i\alpha_i$ and $w_*=\min_iw_i$.  The centers
$y_j=a^j$, $c_{ji}=jr_i$, $0\le j<K$, have nonzero, pairwise distinct
$y_j$, since $a>0$ and $a\ne1$.  This is true when $a<1$ as well.

Choose $F_0>2/\theta$ and define
\begin{equation}
 T_i=\left\lceil F_0w_i/v_0\right\rceil,\qquad
 G_i(t)=\sum_{1\le k<T_i}\frac{(-1)^{k+1}t^k}{k},\qquad
 L_i=\lcm(1,\ldots,T_i-1),\qquad
 D_H=\prod_iL_i^{\lfloor H/w_i\rfloor}.
 \label{rat:truncations}
\end{equation}
An empty sum is zero, and $L_i=1$ when $T_i=1$.
Since $v_0T_i\ge F_0w_i>w_i/\theta$, the substitution
$U_i\mapsto U_i+G_i(t)-\log(1+t)$ and its inverse preserve the
$V$-weighted filtration, where
$V=(v_0,w_1/\theta,\ldots,w_m/\theta)$.
It therefore induces an automorphism on every coefficient packet of
weight less than $H$.  Interpolation remains surjective after replacing
each $\log(1+t)$ by $G_i(t)$.

Use columns $P=Y^hX^\alpha$ with $w_0h+w\alpha\le H$, and rows
$\rho=(j,s,\beta)$ with
$0\le j<K$ and $v_0s+w\beta/\theta<H$.  The entry is
\begin{equation}
 A_{\rho;h,\alpha}
 =[t^sU^\beta]
 P\bigl(a^j(1+t),(jr_i+G_i(t)+U_i)_{i=1}^m\bigr).
 \label{rat:matrix}
\end{equation}
For all sufficiently large, sufficiently divisible $H$, interpolation
supplies a nonzero square minor $\Delta_H$ using every row.  Let $M=M_H$
be its size and put
\begin{equation}
 \bar b=\frac{1}{MH}\sum_{\rho=(j,s,\beta)}w\beta,
 \qquad
 M\sim\frac{K\theta^mH^{m+1}}
 {(m+1)!v_0\prod_iw_i},\qquad 0\le\bar b\le\theta.
 \label{rat:row-count}
\end{equation}
The count follows by the simplex-volume argument used in
\eqref{det:row-count}; the strict row inequality has the same leading
term.  Every limit in $H$ below keeps all the other data fixed.

\subsection{Exact arithmetic clearing for a rational base}

\begin{lemma}[Rational-base arithmetic estimate]
\label{rat:arithmetic}
If $w_0\ge(K-1)\log v$, then
\begin{equation}
 \frac{\log|\Delta_H|}{MH}\ge-(1-\bar b)-E_{\rm ar},
 \qquad
 E_{\rm ar}=\frac{\Lambda F_0m}{v_0}
  +\Lambda\sum_i\frac1{w_i}+\frac{\theta}{w_*},
 \qquad\Lambda=4\log2.
 \label{rat:lower}
\end{equation}
\end{lemma}

\begin{proof}
Scale the column $Y^hX^\alpha$ by
$v^{(K-1)h}\prod_iq_i^{\alpha_i}$ and the row $(j,s,\beta)$ by
$\prod_iq_i^{-\beta_i}$.  Its scaled entry is zero unless
$\alpha\ge\beta$ componentwise; otherwise it is exactly
\begin{equation}
 u^{jh}v^{(K-1-j)h}\prod_i\binom{\alpha_i}{\beta_i}
 [t^s](1+t)^h
 \prod_i\bigl(jp_i+q_iG_i(t)\bigr)^{\alpha_i-\beta_i}.
 \label{rat:scaled-entry}
\end{equation}
Both exponents on $u$ and $v$ are nonnegative.  The denominator of
this entry divides $\prod_iL_i^{\alpha_i-\beta_i}$, which divides
$D_H$.  Multiplying every scaled entry by $D_H$ gives an integer
matrix with nonzero determinant.  Hence
\begin{equation}
\begin{aligned}
 \log|\Delta_H|\ge{}&-M\log D_H
 -\sum_{\text{selected columns }(h,\alpha)}
       \left((K-1)h\log v+\sum_i\alpha_i\log q_i\right)\\
 &+\sum_{\text{rows }(j,s,\beta)}\sum_i\beta_i\log q_i.
\end{aligned}
\label{rat:clearing}
\end{equation}
This is the exact clearing inequality; the numerator $u$ requires no
additional factor.
For every selected column, $\log q_i\le w_i$ and the hypothesis on
$w_0$ give
\[
 (K-1)h\log v+\sum_i\alpha_i\log q_i
 \le w_0h+w\alpha\le H.
\]
Since $0\le w_i-\log q_i<1$ and $w\beta<\theta H$, the row contribution
is at least
\[
 MH\bar b-\sum_{\text{rows}}|\beta|
 \ge MH\bar b-MH\theta/w_*.
\]
The elementary least-common-multiple bound
$\log\lcm(1,\ldots,n)\le\Lambda n$, established in the proof of
Lemma~\ref{det:arithmetic}, gives
\[
 \frac{\log D_H}{H}
 \le\Lambda\sum_i\frac{T_i}{w_i}
 \le\frac{\Lambda F_0m}{v_0}+\Lambda\sum_i\frac1{w_i}.
\]
Substitution in \eqref{rat:clearing} proves \eqref{rat:lower}.
\end{proof}

\subsection{Translation and analytic collision with either sign}

For a transverse multi-index $\gamma$, let
\[
 f_{\gamma,P}(z)=[U^\gamma]P(e^z,z+U_1,\ldots,z+U_m).
\]
These entire functions have entries
$\binom\alpha\gamma e^{hz}z^{|\alpha|-|\gamma|}$ when
$\gamma\le\alpha$, and zero otherwise.
Put $\epsilon_i=j(r_i-\omega)$, $\tau_i=G_i-\log(1+t)$ and
$z=j\omega+\log(1+t)$.  Because $e^{j\omega}=a^j$, direct expansion
gives the row-vector identity
\begin{equation}
\begin{aligned}
 \operatorname{row}_{j,s,\beta}
 ={}&\sum_{\substack{\gamma\ge\beta\\w\gamma\le H}}
       \binom\gamma\beta
       \sum_{0\le d\le\gamma-\beta}\binom{\gamma-\beta}{d}
       \epsilon^{\gamma-\beta-d}
       \sum_{k=0}^s
       \left([t^k]\prod_i\tau_i(t)^{d_i}\right)\\
 &\hspace{23mm}\cdot
       \bigl([t^{s-k}]f_{\gamma,P}(j\omega+\log(1+t))\bigr)_P.
\end{aligned}
\label{rat:row-identity}
\end{equation}
All its scalar coefficients are independent of $P$.  The identity has
at most
\[
 Q_H=(\lfloor H\rfloor+1)^{2m}(\lfloor H/v_0\rfloor+1)
\]
terms, because $w_i\ge1$.  Each nonzero tail coefficient has
$\sum_id_iT_i\le k\le s$, hence
$\sum_id_iw_i<H/F_0$.  On $|t|=1/2$, $|\tau_i(t)|\le1$;
Cauchy's estimate bounds its product's $t^k$ coefficient by $2^k$.
Furthermore
$|\epsilon_i|\le2Ke^\nu e^{-\nu w_i}$ and
$\binom\gamma\beta\binom{\gamma-\beta}{d}\le4^{|\gamma|}$.
Consequently every scalar $\xi$ with index $\gamma$ satisfies
\begin{equation}
 |\xi|\le\exp\{-\nu w(\gamma-\beta)+H E_{\rm tr}\},
 \qquad
 E_{\rm tr}=\frac{\nu}{F_0}+\frac{\log2}{v_0}
       +\frac{\log4+\log(2K)+\nu}{w_*}.
 \label{rat:translation-bound}
\end{equation}
The argument uses absolute values throughout, so it does not require
$\omega>0$ or any sign condition on the $p_i$.

Set
\begin{equation}
 R_a=4K(1+|\log a|),\qquad c=\frac{\log2}{4},\qquad
 E_{{\rm hol},a}=\frac{R_a}{w_0}+\frac{\log2}{v_0}
                   +\frac{\log(2R_a)}{w_*}.
 \label{rat:hol-error}
\end{equation}
Choose one test row from each row expansion, remove its scalar, and
write $n_\gamma$ for the number of resulting rows with index $\gamma$.
For the resulting matrix $T$, we claim
\begin{equation}
 |\det T|\le
 \exp\left\{-c\sum_\gamma n_\gamma^2
                 +MH(E_{{\rm hol},a}+\rho_H)\right\},
 \qquad\rho_H\longrightarrow0,
 \label{rat:collision}
\end{equation}
uniformly over its row choices and the selected columns.
Indeed, on $|z|\le R_a$ the monomial formula gives
\[
 |f_{\gamma,P}(z)|\le
 \exp\left\{H\left(\frac{R_a}{w_0}
                       +\frac{\log(2R_a)}{w_*}\right)\right\}
 =:\mathcal D_H.
\]
Write $f_{\gamma,P}(z)=\sum_{d\ge0}c_{\gamma,d,P}(z/R_a)^d$;
then $|c_{\gamma,d,P}|\le\mathcal D_H$.
The logarithm in the test is its analytic germ on $|t|<1$.
Its power series gives $|\log(1+t)|\le\log2$ on $|t|=1/2$, so
\[
 |j\omega+\log(1+t)|
 \le(K-1)|\omega|+\log2<R_a/2.
\]
Thus a test of order $\ell\le s$ on $(z/R_a)^d$ has modulus at most
$2^\ell2^{-d}$.  The determinant expansion in these series is absolutely
convergent, since each coefficient determinant is bounded by
$M!\mathcal D_H^M$ and the geometric series converge.
Any summand with two equal pairs $(\gamma,d)$ vanishes, since its two
coefficient rows coincide.  Every remaining summand therefore has
$\sum_{\text{rows}}d\ge\sum_\gamma\binom{n_\gamma}{2}$.
Use half of the factor $2^{-\sum d}$ for this saving and sum the other
half freely.  Since $\sum\ell\le MH/v_0$, this yields
\[
 |\det T|\le M!\mathcal D_H^M2^{MH/v_0}
  2^{-\frac12\sum_\gamma\binom{n_\gamma}{2}}
  (1-2^{-1/2})^{-M}.
\]
Equation~\eqref{rat:collision} follows, for example with
\[
 \rho_H=\frac{\log M+\frac14\log2
                       -\log(1-2^{-1/2})}{H}\longrightarrow0,
\]
where \eqref{rat:row-count} gives the limit.

For $0<A<1$ and $0<\eta<1$, put
\begin{equation}
 g=\nu\bigl(A(1-\eta)-\theta\bigr)-(1-\theta),\qquad
 L=\frac{\eta^2K\theta^m}{(m+1)v_0A^m},\qquad
 E_a=E_{\rm ar}+E_{\rm tr}+E_{{\rm hol},a}.
 \label{rat:budgets}
\end{equation}
The contradiction conditions are exactly
\begin{equation}
 w_0\ge(K-1)\log v,\qquad g>0,\qquad E_a<g,
 \qquad cL>1+E_a.
 \label{rat:conditions}
\end{equation}
Here is the full comparison.  Define
$D_A(H)=\#\{\gamma\in\NN^m:w\gamma\le AH\}$ and
$L_H=\eta^2M/(HD_A(H))$.  Simplex counts give
$D_A(H)\sim(AH)^m/(m!\prod_iw_i)$ and $L_H\to L$.
In one term of the multilinear expansion of $\Delta_H$, if at least
$\eta M$ indices have $w\gamma\le AH$, then
$\sum n_\gamma^2\ge(\eta M)^2/D_A(H)=MH L_H$.
The scalar main exponent $-\nu\sum w(\gamma-\beta)$ is nonpositive,
so this term has normalized logarithmic upper bound
$-cL_H+E_{\rm tr}+E_{{\rm hol},a}+\rho_H$.
Otherwise more than $(1-\eta)M$ indices have $w\gamma>AH$.
The scalar estimate then gives normalized main exponent at most
$-\nu(A(1-\eta)-\bar b)$; discard the nonpositive collision term.
There are at most $Q_H^M$ expansion terms.  Taking absolute values
before summing gives
\begin{equation}
 \frac{\log|\Delta_H|}{MH}
 \le E_{\rm tr}+E_{{\rm hol},a}+\varepsilon_H
    +\max\{-cL_H,-\nu(A(1-\eta)-\bar b)\},
 \quad\varepsilon_H=\rho_H+\frac{\log Q_H}{H}\to0.
 \label{rat:upper}
\end{equation}
For the second alternative the gap over $1-\bar b$ is at least $g$,
because
\[
 \nu(A(1-\eta)-\bar b)-(1-\bar b)
 =\nu A(1-\eta)-1-(\nu-1)\bar b\ge g.
\]
For the first alternative use $1-\bar b\le1$ and $L_H\to L$.
Under \eqref{rat:conditions}, both upper bounds in
\eqref{rat:upper} are strictly less than the lower bound
$-(1-\bar b)-E_{\rm ar}$ for sufficiently large $H$.  This
contradicts \eqref{rat:lower}.

\subsection{Finite parameter choice for each fixed base}

\begin{proof}[Proof of Theorem~\ref{rat:main}]
Fix $a=u/v\in\QQ_{>0}\setminus\{1\}$ and $\nu>2$ before making any
dimension or approximation choice.  Suppose that arbitrarily large
integers $q$ admit $p\in\ZZ$ with $|\log a-p/q|\le q^{-\nu}$.
Choose $\theta,A,B,C,\eta$ as in Lemma~\ref{par:constants}, so
\[
 0<\theta<A<B<1,\qquad C>1,\qquad CB<1,\qquad B<C\theta<1,
 \qquad g>0.
\]
Let $\varepsilon_0=\min(g/2,1/2)$ and choose $F_0>2/\theta$ such that
$\nu/F_0<\varepsilon_0/3$.
For a variable positive integer $m$, set
\[
 K=\lfloor C^m\rfloor,\qquad w_0=B^{-m},\qquad
 v_0=2K\theta^m w_0.
\]
The interpolation volume hypotheses hold, since
$K\theta^m\le(C\theta)^m<1$ and
$K(w_0/v_0)\theta^m=1/2$.
Moreover
\[
 \frac{K}{w_0}\le(CB)^m\to0,\qquad
 v_0\sim2(C\theta/B)^m\to\infty,\qquad m/v_0\to0,
 \qquad
 L=\frac{\eta^2(B/A)^m}{2(m+1)}\to\infty.
\]
Because $a$ was already fixed, these limits also give
\[
 R_a/w_0=4(1+|\log a|)K/w_0\to0,\qquad
 (K-1)\log v/w_0\to0.
\]
Consequently one finite $m$ can be chosen such that
\begin{equation}
 w_0\ge(K-1)\log v,\qquad
 \frac{\Lambda F_0m+2\log2}{v_0}+\frac{R_a}{w_0}
       <\varepsilon_0/3,\qquad cL>2.
 \label{rat:dimension-margin}
\end{equation}
Fix this $m$, hence $K,w_0,v_0,R_a$ and all the constants in this display.

Choose $m$ approximations as in \eqref{rat:approximations}
successively, with denominators sufficiently large to satisfy the
successive weight thresholds of Theorem~\ref{geom:interpolation} and
\begin{equation}
 \Lambda\sum_i\frac1{w_i}
 +\frac{\theta+\log4+\log(2K)+\nu+\log(2R_a)}{w_*}
       <\varepsilon_0/3.
 \label{rat:weight-margin}
\end{equation}
To realize the last inequality, first impose a sufficiently large common
lower bound on all $w_i$, using $\sum_i1/w_i\le m/w_*$.
Then impose the geometric threshold for each weight after its predecessors
have been chosen.  The assumed good denominators are unbounded, so every
one of these finitely many lower bounds is attainable, even if those
denominators are sparse.  The geometric weight thresholds are independent
of the centers subsequently formed from these choices.
No nonzero-numerator hypothesis is needed: the distinctness of $a^j$
does not involve the $p_i$.

The total error is exactly
\[
\begin{aligned}
 E_a={}&\frac{\nu}{F_0}
  +\frac{\Lambda F_0m+2\log2}{v_0}+\frac{R_a}{w_0}
  +\Lambda\sum_i\frac1{w_i}\\
 &+\frac{\theta+\log4+\log(2K)+\nu+\log(2R_a)}{w_*}
 <\varepsilon_0.
\end{aligned}
\]
Thus $E_a<g$ and $cL>2>1+E_a$, and all conditions
\eqref{rat:conditions} hold.  After fixing these weights and centers,
let $H$ tend to infinity through sufficiently divisible integers above
the interpolation threshold for these particular centers.  The comparison
\eqref{rat:upper}--\eqref{rat:lower} is a contradiction.  No uniform
threshold in $H$ over different bases or centers has been used.

The good denominators for the fixed $a,\nu$ are therefore bounded,
which proves \eqref{rat:eventual-bound}.  If $\omega$ were rational,
its exact representations with arbitrarily large positive denominators
would contradict that bound, so $\omega$ is irrational.
It follows that $\mu(\omega)\le2$.
For completeness, pigeonhole applied to $0,\omega,\ldots,N\omega$
modulo one gives $p\in\ZZ$ and $1\le q\le N$ with
$0<|\omega-p/q|\le1/(qN)\le1/q^2$.
The errors tend to zero as $N\to\infty$; after reduction the denominators
must be unbounded, since $\omega$ is irrational.
Every exponent less than two consequently occurs infinitely often,
so $\mu(\omega)\ge2$.  This completes the conditional proof.
\end{proof}

\subsection{Rational scaling and further conditional consequences}

\begin{lemma}[Invariance under nonzero rational scaling]
\label{rat:scaling}
If $x$ is an irrational real number and $b\in\QQ\setminus\{0\}$,
then $bx$ is irrational and $\mu(bx)=\mu(x)$.
\end{lemma}

\begin{proof}
Write $b=r/s$ with $r\in\ZZ\setminus\{0\}$, $s\ge1$ and
$\gcd(|r|,s)=1$.  Irrationality of $bx$ follows by division by $b$.
Suppose a positive exponent $\tau$ occurs infinitely often for $x$,
with coprime $p,q$, $q\ge2$, and
$0<|x-p/q|<q^{-\tau}$.
Reduce $rp/(sq)$ to $P/Q$, and set $d=\gcd(|rp|,sq)$.
Since $\gcd(p,q)=1$, a Bezout identity for $p,q$ shows that $d$ divides
$rs$: multiply that identity by $rs$, noting that $d$ divides both
$rp$ and $sq$.  Thus $1\le d\le|r|s$ and
\[
 q/|r|\le Q=sq/d\le sq.
\]
In particular $Q$ is unbounded.  Regardless of the sign of $r$,
\[
 0<\left|bx-\frac PQ\right|
 =\frac{|r|}{s}\left|x-\frac pq\right|
 <\frac{|r|}{s}q^{-\tau}
 \le |r|s^{\tau-1}Q^{-\tau}.
\]
For every $0<\lambda<\tau$, the fixed constant $|r|s^{\tau-1}$ is
less than $Q^{\tau-\lambda}$ for all sufficiently large $Q$.
Hence $\lambda$ occurs infinitely often for $bx$.
Taking suprema and then letting $\lambda$ increase to each such
$\tau$ gives $\mu(bx)\ge\mu(x)$, also when that supremum is infinite.
Apply the same argument to $1/b$ and $bx$ to obtain the reverse inequality.
\end{proof}

\begin{corollary}[Conditional consequences; candidate]
\label{rat:span}
Assume Theorem~\ref{geom:interpolation} with its stated quantifiers.
Every nonzero finite rational linear combination
\[
 \omega=\sum_{k=1}^n b_k\log a_k,
 \qquad b_k\in\QQ,\quad a_k\in\QQ_{>0},
\]
is irrational and has $\mu(\omega)=2$.
In particular this applies to $\log a$ whenever $a>0$ is algebraic,
$a\ne1$, and $a^N\in\QQ$ for some integer $N\ge1$.
\end{corollary}

\begin{proof}
Choose a positive integer $N$ such that all $Nb_k$ are integers.
The real logarithm identities give
\[
 N\omega=\log r,\qquad
 r=\prod_{k=1}^n a_k^{Nb_k}\in\QQ_{>0}.
\]
Negative integer exponents preserve positive rationality.  Since
$\omega\ne0$, one has $r\ne1$.  Theorem~\ref{rat:main} gives
irrationality and exponent two for $\log r$, conditionally on the
interpolation hypothesis.  Lemma~\ref{rat:scaling} with the factor
$1/N$ gives the assertion for $\omega$.
For the final statement, $r=a^N$ is positive rational and is not one,
and $\log a=(1/N)\log r$; apply the same argument.
\end{proof}

The conditional statements above add no interpolation dependency beyond
Theorem~\ref{geom:interpolation}.  They assert a base-dependent eventual
denominator threshold, rather than a common threshold over all rational
bases.  They make no claim about logarithms of arbitrary positive
algebraic numbers whose positive integer powers are never rational.


\section{Audit interface for the candidate proof}
The key dependency is Theorem~\ref{geom:interpolation}, including its
center-independent successive weight thresholds. Its proposed proof
passes through the local transverse multiplicity bound, persistent
derivative-ideal components, the residue obstruction forcing constant
$Y$, and the ordinary blow-up and Serre-vanishing interpolation argument.
Replacing $Y=1$ by $Y=2^j$ is a new mathematical step.

If that geometric statement and the determinant estimates hold as
written, Section 4 selects a finite dimension for each fixed $\nu>2$,
then chooses approximation denominators and their separated weights,
then fixes the centers, and only finally lets $H$ tend to infinity.
Excluding unbounded $\nu$-good approximations for every $\nu>2$ gives
$\mu(\log2)\le2$; the final Dirichlet argument gives the reverse bound.
The endpoint equality is therefore a conditional consequence of the
complete proposed chain, not a deduction from finite numerical tests.

Internal readers have not located a definite invalid inference in the
inspected chain. That observation is not proof certification. The new
interpolation theorem, its hypotheses, and its use in the determinant
contradiction should be checked before citing this document as a result.

\clearpage
\begin{thebibliography}{99}
\bibitem{openai} OpenAI, \emph{The irrationality exponent of $\pi$ is $2$},
24 September 2026. \href{https://github.com/openai/math/blob/main/preprints/The-irrationality-exponent-of-pi-is-2-September-24-2026/paper.pdf}{Manuscript and source repository}.
\bibitem{marcovecchio} R. Marcovecchio, \emph{The Rhin--Viola method for log 2},
Acta Arith. 139 (2009), 147--184.
\href{https://www.impan.pl/shop/en/publication/transaction/download/product/82234}{Publisher PDF}.
\bibitem{mondal2021} P. Mondal, \emph{How many zeroes?}, Draft III,
\href{https://arxiv.org/abs/1806.05346v5}{arXiv:1806.05346v5},
Corollary VIII.3, for isolated local intersection multiplicities.
\bibitem{stacks} The Stacks Project,
\href{https://stacks.math.columbia.edu/tag/02OS}{02OS},
\href{https://stacks.math.columbia.edu/tag/02ND}{02ND},
\href{https://stacks.math.columbia.edu/tag/02NS}{02NS},
\href{https://stacks.math.columbia.edu/tag/0AG6}{0AG6},
\href{https://stacks.math.columbia.edu/tag/0AG7}{0AG7},
\href{https://stacks.math.columbia.edu/tag/0B5U}{0B5U},
\href{https://stacks.math.columbia.edu/tag/0BXX}{0BXX}.
\bibitem{lazarsfeld} R. Lazarsfeld, \emph{Positivity in Algebraic Geometry I},
Springer, 2004, Corollary 1.4.10.
\bibitem{demailly1992} J.-P. Demailly, \emph{Singular hermitian metrics on positive
line bundles}, LNM 1507 (1992), 87--104, Section 6.
\bibitem{philippon1986} P. Philippon, \emph{Lemmes de z\'eros dans les groupes
alg\'ebriques commutatifs}, Bull. SMF 114 (1986), 355--383; errata 115 (1987).
\bibitem{laurent1991} M. Laurent, \emph{Sur quelques r\'esultats r\'ecents de
transcendance}, Ast\'erisque 198--200 (1991), 209--230, Section 6.
\bibitem{laurent2000} M. Laurent, \emph{Interpolation determinants of exponential
polynomials}, Publ. Math. Debrecen 56 (2000), 457--473, Theorem 1.
\end{thebibliography}
\end{document}
